\documentclass[10pt]{article}
\usepackage[T1]{fontenc}
\usepackage{lmodern}
\usepackage[margin=22mm]{geometry}
\usepackage{xcolor}
\usepackage{microtype}
\usepackage{enumitem}
\usepackage{hyperref}
\usepackage{xurl}
\usepackage{forest}
\usetikzlibrary{arrows.meta}
\usepackage{pdflscape}
\usepackage[most]{tcolorbox}
\usepackage{titlesec}
\usepackage{fancyhdr}
\definecolor{Navy}{HTML}{17324D}
\definecolor{Teal}{HTML}{176B6B}
\definecolor{PaleBlue}{HTML}{EAF2F8}
\definecolor{FileFill}{HTML}{F7F9FA}
\definecolor{DescText}{HTML}{334E5C}
\definecolor{TreeLine}{HTML}{607D8B}
\definecolor{Warn}{HTML}{8A4B08}
\hypersetup{colorlinks=true,linkcolor=Teal,urlcolor=Teal,pdfborder={0 0 0}}
\titleformat{\section}{\Large\bfseries\color{Navy}}{}{0pt}{}[\color{Teal}\titlerule]
\titleformat{\subsection}{\large\bfseries\color{Teal}}{}{0pt}{}
\setlist[itemize]{leftmargin=*,itemsep=.45em,topsep=.45em}
\setlist[enumerate]{leftmargin=*,itemsep=.45em,topsep=.45em}
\setlength{\parindent}{0pt}
\setlength{\parskip}{.65em}
\newcommand{\code}[1]{\texttt{#1}}
\forestset{
 rootnode/.style={draw=Navy,fill=Navy,text=white,rounded corners=2pt,inner sep=5pt},
 filenode/.style={draw=TreeLine,fill=FileFill,rounded corners=2pt,inner sep=3.2pt},
}
\pagestyle{fancy}\fancyhf{}\fancyhead[L]{\small\color{DescText}Repository Structure Inventory}\fancyhead[R]{\small\color{DescText}\leftmark}\fancyfoot[C]{\thepage}
\setlength{\headheight}{14pt}
\newtcolorbox{verificationbox}{breakable,colback=orange!5,colframe=Warn,title={Verification status},fonttitle=\bfseries}
\begin{document}
\hypersetup{pageanchor=false}
\begin{titlepage}
\pagecolor{PaleBlue}
\color{Navy}
\vspace*{0.17\textheight}
{\Huge\bfseries Grade 10 — Term 1 Structure\par}
\vspace{1em}
{\Large Repository resource inventory\par}
\vfill
{\large Authoritative conversion of \code{STRUCTURE.md}\par}
\end{titlepage}
\hypersetup{pageanchor=true}
\nopagecolor
\tableofcontents
\clearpage

\section{Overview}

\code{term\_1} is a LaTeX source collection for a Grade 10 Global Economics unit on decision analysis. It progresses from C1–C2 practice (identifying decision components and constructing payoff tables), through C1–C10 assessments and worked solutions, to criterion-specific recovery activities and supplementary material. Files described as \textbf{student-facing} are worksheets, assessments, or briefs intended to receive student work; files described as \textbf{teacher material} expose worked answers or instructional solution sets. There are no local HTML pages, stylesheets, JavaScript programs, datasets, configuration files, or standalone media: the JavaScript item is a student project \emph{brief}, not an implementation.

\section{Directory tree}

\begin{tcolorbox}[colback=PaleBlue,colframe=Teal,title={Reading the diagrams}]
The directory tree is divided into one landscape diagram per top-level directory so that every documented entry and description remains legible. Directory nodes use a dark fill; file nodes use a light fill, with descriptions shown in a secondary text color. The inventory root is \code{.}.
\end{tcolorbox}
\begin{landscape}
\subsection{\code{1-practice\_activities/}}
\begin{center}
\begin{forest}
for tree={grow=east, parent anchor=east, child anchor=west, anchor=west, edge={-Latex,draw=TreeLine}, l sep=9mm, s sep=1.2mm, font=\sffamily\fontsize{6.2}{7.2}\selectfont},
[{\textbf{1-practice\_activities/}\\\textcolor{DescText}{Student practice and an instructor-facing contents summary for Lesson 1 decision components/payoff tables.}}, rootnode, text width=6.0cm
[{\code{G10\_T1\_L1\_C1-C2\_practice\_activity1.tex}\\\textcolor{DescText}{Student-facing C1–C2 workbook: 21 scenarios spanning direct, variable, fixed, and hybrid payoff models, with guided interpretations and computed tables.}}, filenode, text width=16cm]
[{\code{G10\_T1\_L1\_C1-C2\_practice\_activity2 - Copy.tex}\\\textcolor{DescText}{Compact alternate/legacy Activity 2: 14 terse scenario notes and partial calculations based largely on Activity 1 contexts; unlike the full Activity 2, it is not a nine-problem worked worksheet.}}, filenode, text width=16cm]
[{\code{G10\_T1\_L1\_C1-C2\_practice\_activity2.tex}\\\textcolor{DescText}{Student-facing worked C1–C2 activity with nine new contexts (food truck through emergency supplies), each identifying decision elements and building a payoff table.}}, filenode, text width=16cm]
[{\code{G10\_T1\_summary.tex}\\\textcolor{DescText}{Teacher/contributor summary that inventories the modeling categories and scenario counts in the two Lesson 1 practice activities.}}, filenode, text width=16cm]
]
\end{forest}
\end{center}
\end{landscape}
\begin{landscape}
\subsection{\code{2-learning\_evidences/}}
\begin{center}
\begin{forest}
for tree={grow=east, parent anchor=east, child anchor=west, anchor=west, edge={-Latex,draw=TreeLine}, l sep=9mm, s sep=1.2mm, font=\sffamily\fontsize{6.2}{7.2}\selectfont},
[{\textbf{2-learning\_evidences/}\\\textcolor{DescText}{Student assessments/project briefs and paired teacher solution documents for Lessons 1–3.}}, rootnode, text width=6.0cm
[{\code{G10\_T1\_L1\_exam.tex}\\\textcolor{DescText}{Student-facing Learning Evidence 1 exam: three decision problems assessing C1–C5 via payoff tables, expected value, maximax, and maximin.}}, filenode, text width=16cm]
[{\code{G10\_T1\_L1\_exam\_solution.tex}\\\textcolor{DescText}{Teacher material with full C1–C5 workings for the exam’s cafeteria, delivery-cooperative, and sports-center problems.}}, filenode, text width=16cm]
[{\code{G10\_T1\_L1\_practice\_exam.tex}\\\textcolor{DescText}{Student-facing two-problem C1–C5 practice exam on recreation-center memberships and fundraiser menu plans.}}, filenode, text width=16cm]
[{\code{G10\_T1\_L1\_practice\_exam\_solution.tex}\\\textcolor{DescText}{Teacher material with complete decision-component, payoff, maximax, maximin, and expected-value solutions to that practice exam.}}, filenode, text width=16cm]
[{\code{G10\_T1\_L2\_C6\_activity.tex}\\\textcolor{DescText}{Student-facing four-problem C6 activity applying maximum opportunity/minimax regret to bakery, farm-storage, theater, and another payoff-table case.}}, filenode, text width=16cm]
[{\code{G10\_T1\_L2\_C6\_activity\_solution.tex}\\\textcolor{DescText}{Teacher material giving full and abbreviated minimax-regret solutions for the four C6 problems.}}, filenode, text width=16cm]
[{\code{G10\_T1\_L2\_C7\_activity.tex}\\\textcolor{DescText}{Student-facing ten-problem C7 expected-value activity, from one-project choices to multi-alternative and joint-state decisions.}}, filenode, text width=16cm]
[{\code{G10\_T1\_L2\_C7\_activity\_solution.tex}\\\textcolor{DescText}{Teacher material with worked expected-value analysis for all ten C7 activity problems.}}, filenode, text width=16cm]
[{\code{G10\_T1\_L2\_midterm.tex}\\\textcolor{DescText}{Student-facing Lesson 2 midterm assessing C1–C7 across payoff construction, classical criteria, regret, and expected-value problems.}}, filenode, text width=16cm]
[{\code{G10\_T1\_L2\_midterm\_solution.tex}\\\textcolor{DescText}{Teacher material with worked solutions to eight midterm scenarios, including bakery, shipping, theater, inventory, routing, marketing, production, and zoning choices.}}, filenode, text width=16cm]
[{\code{G10\_T1\_L2\_mock\_midterm.tex}\\\textcolor{DescText}{Student-facing mock midterm for C1–C7; uses a smaller alternate set of logistics, inventory, routing, production, and zoning problems.}}, filenode, text width=16cm]
[{\code{G10\_T1\_L2\_mock\_midterm\_solution.tex}\\\textcolor{DescText}{Teacher material with full solutions for the mock midterm’s five scenarios.}}, filenode, text width=16cm]
[{\code{G10\_T1\_L3\_proyect.tex}\\\textcolor{DescText}{Student-facing project brief requiring a JavaScript decision-analysis tool, presentation, and decision-tree defense for C8–C10; includes specifications and solved examples, but no JavaScript source.}}, filenode, text width=16cm]
[{\code{G10\_T1\_L3\_proyect\_problems.tex}\\\textcolor{DescText}{Student-facing bank of ten project datasets/tasks for comparing Bayes, Maximin, and Minimax Regret and constructing decision trees.}}, filenode, text width=16cm]
[{\code{G10\_T1\_L3\_proyect\_problems\_solutions.tex}\\\textcolor{DescText}{Teacher material with extensive worked C8/C9 solutions and C10 evidence guidance for all ten project problems.}}, filenode, text width=16cm]
[{\code{G10\_T1\_L3\_proyect\_problems\_solutions\_min.tex}\\\textcolor{DescText}{Teacher-facing condensed solution variant: shorter calculations/recommendations for the same ten project problems rather than the full explanatory treatment.}}, filenode, text width=16cm]
]
\end{forest}
\end{center}
\end{landscape}
\begin{landscape}
\subsection{\code{3-catch\_ups/}}
\begin{center}
\begin{forest}
for tree={grow=east, parent anchor=east, child anchor=west, anchor=west, edge={-Latex,draw=TreeLine}, l sep=9mm, s sep=1.2mm, font=\sffamily\fontsize{6.2}{7.2}\selectfont},
[{\textbf{3-catch\_ups/}\\\textcolor{DescText}{Criterion-by-criterion recovery worksheets (student-facing) and their teacher answer keys for C1–C7.}}, rootnode, text width=6.0cm
[{\code{G10\_T1\_C1\_catchUp.tex}\\\textcolor{DescText}{Student-facing ten-problem C1 recovery activity on formulating choices under uncertainty and using expected value.}}, filenode, text width=16cm]
[{\code{G10\_T1\_C1\_catchUp\_soution.tex}\\\textcolor{DescText}{Teacher answer key for C1, with expected-value computations and recommendations (\code{soution} is the stored filename).}}, filenode, text width=16cm]
[{\code{G10\_T1\_C2\_catchUp.tex}\\\textcolor{DescText}{Student-facing ten-problem C2 recovery activity identifying alternatives, events, consequences, and states of nature.}}, filenode, text width=16cm]
[{\code{G10\_T1\_C2\_catchUp\_soution.tex}\\\textcolor{DescText}{Teacher answer key supplying the C2 decision-component interpretations (\code{soution} is the stored filename).}}, filenode, text width=16cm]
[{\code{G10\_T1\_C3\_catchUp.tex}\\\textcolor{DescText}{Student-facing ten-problem C3 recovery activity for turning narrative and supplied data into payoff tables.}}, filenode, text width=16cm]
[{\code{G10\_T1\_C3\_catchUp\_soution.tex}\\\textcolor{DescText}{Teacher answer key with constructed payoff tables and supporting calculations for C3 (\code{soution} is the stored filename).}}, filenode, text width=16cm]
[{\code{G10\_T1\_C4\_catchUp.tex}\\\textcolor{DescText}{Student-facing ten-problem C4 recovery worksheet applying the optimistic maximax rule without probabilities.}}, filenode, text width=16cm]
[{\code{G10\_T1\_C4\_catchUp\_soution.tex}\\\textcolor{DescText}{Teacher key showing each row maximum and maximax choice for C4 (\code{soution} is the stored filename).}}, filenode, text width=16cm]
[{\code{G10\_T1\_C5\_catchUp.tex}\\\textcolor{DescText}{Student-facing ten-problem C5 recovery worksheet applying the conservative maximin rule without probabilities.}}, filenode, text width=16cm]
[{\code{G10\_T1\_C5\_catchUp\_solution.tex}\\\textcolor{DescText}{Teacher key showing each row minimum and maximin recommendation for C5.}}, filenode, text width=16cm]
[{\code{G10\_T1\_C6\_catchUp.tex}\\\textcolor{DescText}{Student-facing ten-problem C6 recovery worksheet using maximum opportunity/minimax regret.}}, filenode, text width=16cm]
[{\code{G10\_T1\_C6\_catchUp\_soution.tex}\\\textcolor{DescText}{Teacher key with regret tables and minimum-of-maximum-regret choices for C6 (\code{soution} is the stored filename).}}, filenode, text width=16cm]
[{\code{G10\_T1\_C7\_catchUp.tex}\\\textcolor{DescText}{Student-facing ten-problem C7 recovery activity comparing alternatives with probabilities and expected value.}}, filenode, text width=16cm]
[{\code{G10\_T1\_C7\_catchUp\_solution.tex}\\\textcolor{DescText}{Teacher key with worked expected values for the ten C7 scenarios.}}, filenode, text width=16cm]
]
\end{forest}
\end{center}
\end{landscape}
\begin{landscape}
\subsection{\code{4-extra\_material/}}
\begin{center}
\begin{forest}
for tree={grow=east, parent anchor=east, child anchor=west, anchor=west, edge={-Latex,draw=TreeLine}, l sep=9mm, s sep=1.2mm, font=\sffamily\fontsize{6.2}{7.2}\selectfont},
[{\textbf{4-extra\_material/}\\\textcolor{DescText}{Supplementary exercises, alternate/reference answers, and worked enrichment material for Lessons 1–3.}}, rootnode, text width=6.0cm
[{\code{G10\_T1\_L1\_V2\_solution.tex}\\\textcolor{DescText}{Teacher material: alternate C1–C5 exam solution version with changed state sets/probabilities and expanded alternatives relative to the main Lesson 1 exam key.}}, filenode, text width=16cm]
[{\code{G10\_T1\_L1\_supplementary.tex}\\\textcolor{DescText}{Student-facing five-problem C1–C5 supplementary assessment (smoothie, shirts, streaming, courier, and water-utility decisions).}}, filenode, text width=16cm]
[{\code{G10\_T1\_L1\_supplementary\_hard\_solution.tex}\\\textcolor{DescText}{Teacher material for a harder parallel five-problem set (coffee cart, bookstore, meal subscriptions, and more complex joint-state cases), not the direct key to the student supplementary sheet.}}, filenode, text width=16cm]
[{\code{G10\_T1\_L1\_supplementary\_solution.tex}\\\textcolor{DescText}{Teacher key matching the standard supplementary contexts and providing full C1–C5 calculations.}}, filenode, text width=16cm]
[{\code{G10\_T1\_L2\_C6\_activity\_supplementary.tex}\\\textcolor{DescText}{Student-facing seven-problem supplementary C6 assessment using increasingly large minimax-regret payoff tables.}}, filenode, text width=16cm]
[{\code{G10\_T1\_L2\_C6\_activity\_supplementary\_solution.tex}\\\textcolor{DescText}{Teacher key with state maxima, regret tables, and minimax-regret choices for those seven C6 problems.}}, filenode, text width=16cm]
[{\code{G10\_T1\_L2\_MATERIAL.tex}\\\textcolor{DescText}{Teacher/instructional worked set combining four C6 minimax-regret examples and ten C7 expected-value examples, with cross-criterion comparisons.}}, filenode, text width=16cm]
[{\code{G10\_T1\_L2\_reference\_midterm.tex}\\\textcolor{DescText}{Teacher reference solution set for C1–C7, organized around café, cafeteria, logistics, investment, production, bookstore, and related decisions; it is not a blank student midterm.}}, filenode, text width=16cm]
[{\code{G10\_T1\_L3\_MATERIAL.tex}\\\textcolor{DescText}{Teacher/instructional worked Lesson 3 material for C8 method comparison, C9 staged decision trees, and C10 end-to-end analysis.}}, filenode, text width=16cm]
]
\end{forest}
\end{center}
\end{landscape}

\section{Important relationships}

\begin{itemize}
\item All entries are standalone LaTeX documents. Solution files are teacher companions to the like-named student assessment/activity unless a different relationship is explicitly noted above; none is loaded by another file in this directory.
\item Most solution, supplementary, and material documents load the shared \code{preamble/preamble\_exams.tex} located outside \code{term\_1}. Catch-up keys reach it with \code{../../../preamble/preamble\_exams.tex}; most other keys/materials use \code{../../preamble/preamble\_exams.tex}.
\item The practice activities and summary load the external shared \code{preamble.tex} through a multi-directory \code{\textbackslash{}input@path}. Activity 2, its compact copy, and the summary additionally load three shared graph-macro files: \path{graphs/uniform_linear_probability_distributions.tex}, \path{graphs/probability_distribution_graphs.tex}, and \path{graphs/normal_distribution_graphs.tex}. They define reusable graph commands; no graph image files are stored here.
\item Student assessment sheets that define their own page layout embed the shared \code{preamble/logo.png} via relative paths. There are no images or other media physically inside \code{term\_1}.
\item \code{G10\_T1\_L3\_proyect.tex} specifies interactive web-page elements and JavaScript functions as student deliverables. The directory contains no HTML, CSS, or JavaScript implementation to which the brief links.
\item \code{G10\_T1\_L3\_proyect\_problems.tex} supplies the ten datasets/tasks used by the full and minimal project solution documents. \code{G10\_T1\_L3\_MATERIAL.tex} is separate worked enrichment rather than a generated build artifact.
\item No generated LaTeX outputs (\code{.aux}, \code{.log}, \code{.out}, \code{.toc}, or PDFs), dependency directories, editor caches, or meaningful hidden files are present. \code{STRUCTURE.md} is intentionally omitted from the inventory tree because it is this inventory document.
\end{itemize}

\section{Classification summary}

\begin{itemize}
\item \textbf{Student-facing pages/documents:} blank activities, exams, midterms, catch-ups, supplementary worksheets, and project briefs/problem banks identified above.
\item \textbf{Teacher materials:} all worked \code{solution}/\code{soution} files, reference/material sets, and the practice summary.
\item \textbf{Stylesheets / JavaScript / data / configuration / local media:} none present. LaTeX preambles, graph definitions, and the logo are shared dependencies outside this directory.
\item \textbf{Generated or build-related files:} none present.
\end{itemize}
\end{document}
